\documentclass[aspectratio=169]{beamer}
\usepackage{booktabs}
\usepackage{hyperref}

\title{Cerebro SDK Ops Briefing}
\subtitle{Operational Playbook for November 6, 2025}
\author{Writer Security Operations}
\date{Updated November 5, 2025}

\begin{document}

\begin{frame}
  \titlepage
\end{frame}

\begin{frame}{Agenda}
  \begin{itemize}
    \item Why the Cerebro SDK matters for Security Ops
    \item Where it fits in daily workflows and incident response
    \item Key surfaces to monitor and the guardrails in place
    \item Rollout plan, support channels, and next actions
  \end{itemize}
\end{frame}

\section{Orientation}

\begin{frame}{What the SDK Delivers}
  \begin{itemize}
    \item A unified way to interact with Cerebro agents, findings, tickets, and analytics.
    \item Consistent models eliminate ad-hoc API calls and one-off scripts.
    \item Enables faster playbook automation while keeping controls centralized.
  \end{itemize}
\end{frame}

\begin{frame}{Why Ops Should Care}
  \begin{columns}
    \begin{column}{0.55\textwidth}
      \begin{itemize}
        \item Shorter triage cycles: surface review queue status and SLA risk in one call.
        \item Better handoffs: align tickets, notifications, and task tracking from a single toolkit.
        \item Stronger audit posture: SDK enforces structured logging and error handling.
      \end{itemize}
    \end{column}
    \begin{column}{0.4\textwidth}
      \begin{block}{Ops Wins}
        \begin{itemize}
          \item Reliable runbook automation
          \item Standardized reporting
          \item Confidence in guardrails
        \end{itemize}
      \end{block}
    \end{column}
  \end{columns}
\end{frame}

\section{Ops Workflows}

\begin{frame}{Daily Operations}
  \begin{itemize}
    \item Review queue overview: pull priorities, assignees, and overdue tasks each morning.
    \item Integration coverage: spot connectors trending toward unhealthy before alerts fire.
    \item Health snapshots: fetch runtime analytics to confirm agent services remain steady.
  \end{itemize}
\end{frame}

\begin{frame}{Incident Response Support}
  \begin{itemize}
    \item Retrieve session timelines to reconstruct decision paths during retros.
    \item Export review history for responders without granting database access.
    \item Trigger targeted notifications or ticket updates straight from incident channels.
  \end{itemize}
\end{frame}

\begin{frame}{Cross-Team Collaboration}
  \begin{itemize}
    \item Share consistent status with SecEng, product, and compliance using the same interfaces.
    \item Leverage in-memory adapters for tabletop exercises without touching production.
    \item Use workflow evaluations to preview automation before requesting deployment.
  \end{itemize}
\end{frame}

\section{Controls and Guardrails}

\begin{frame}{Access \& Permissions}
  \begin{itemize}
    \item SDK honors existing roles: no elevated access beyond what tokens permit.
    \item Encourage short-lived tokens sourced from the standard auth pipeline.
    \item CSRF tokens and audit stamps are surfaced for logging into our SIEM.
  \end{itemize}
\end{frame}

\begin{frame}{Change Management}
  \begin{itemize}
    \item Changesets enforce review before SDK releases ship to teams.
    \item Versioned exports mean automation can pin to a known behavior window.
    \item Recommended: align SDK upgrades with weekly ops change window.
  \end{itemize}
\end{frame}

\begin{frame}{Monitoring Signals}
  \begin{itemize}
    \item Runtime health metrics highlight noisy agents, error spikes, or warning trends.
    \item Integration coverage ratios inform partnership outreach and customer comms.
    \item Task status endpoints surface delayed jobs so Ops can intervene early.
  \end{itemize}
\end{frame}

\section{Adoption Plan}

\begin{frame}{Rollout Checklist}
  \begin{itemize}
    \item Identify initial playbooks (review queue digest, ticket sync, incident export).
    \item Pair Ops analysts with SecEng buddies for the first automation build.
    \item Publish quick links in the Ops runbook to the SDK scripts and dashboards.
    \item Schedule retro one week after launch to capture feedback and close gaps.
  \end{itemize}
\end{frame}

\begin{frame}{Support Channels}
  \begin{itemize}
    \item Slack: \texttt{\#cerebro-ops} for questions, runbook reviews, and status updates.
    \item Office hours every Tuesday with the SDK maintainers for pairing and demos.
    \item PagerDuty layer remains unchanged; SDK adds data, not another on-call rotation.
  \end{itemize}
\end{frame}

\section{Next Steps}

\begin{frame}{Action Items}
  \begin{itemize}
    \item Ops leads nominate two candidate workflows by end-of-week.
    \item Request access tokens through the existing secrets process if not already provisioned.
    \item Drop automation ideas or concerns in the shared backlog doc before Friday.
  \end{itemize}
\end{frame}

\begin{frame}{Q\&A}
  \centering
  Reach us at \texttt{sdk@writer.com}
\end{frame}

\end{document}
